\chapter{Methodology, Requirements Analysis, and Design Constraints}

This chapter defines the methodological framework adopted for the project, derives the
functional and non-functional requirements from the operational context established in
Chapter~2, and documents the design constraints that bounded the solution space. Together
these elements provide the formal basis for the architectural decisions described in
Chapter~4.

\section{Methodology options considered}

The project combines software engineering, browser automation, prompt design, and
evidence-oriented investigation. Three methodologies were considered: Agile, TDSP,
and CRISP-DM.

\subsection{Agile}

Agile matched the day-to-day iteration loop (build, run, inspect, improve), especially
during prompt and tool refinement.

\subsection{TDSP}

TDSP was useful for thinking about experimentation and lifecycle discipline, but the
project is not centered on training a predictive model.

\subsection{CRISP-DM}

CRISP-DM~\cite{crispdm} offered the best end-to-end narrative from business objective
to deployment and evaluation, while still allowing iterative engineering inside each phase.

\begin{table}[H]
\centering
\begin{tabularx}{\textwidth}{>{\raggedright\arraybackslash\bfseries}p{2.7cm}
                              >{\raggedright\arraybackslash}X
                              >{\raggedright\arraybackslash}X}
\toprule
\rowcolor{accent!8}
Methodology & Strength for this project & Limitation for this project \\
\midrule
Agile & Excellent short-cycle execution and adaptation & Less explicit about
  data/evidence lifecycle structure \\
TDSP & Strong experimentation and operationalization framing & More model-training
  oriented than this use case \\
CRISP-DM & Clear business-to-deployment and evaluation progression & Needed adaptation
  for agentic browser workflows \\
\bottomrule
\end{tabularx}
\caption{Methodology comparison used in the analysis phase.}
\label{tab:method-compare}
\end{table}

\section{Chosen methodology and adaptation}

CRISP-DM was selected as the primary methodology, with Agile iteration behavior inside
each phase. The adapted cycle is illustrated in Figure~\ref{fig:crispdm}.

\begin{figure}[H]
\centering
\begin{tikzpicture}[
  phase/.style={draw=accent, fill=blue!8, rounded corners=6pt,
                minimum width=3.5cm, minimum height=1.1cm,
                text centered, font=\small\bfseries, align=center},
  annot/.style={font=\tiny, text=gray!80, align=center, text width=2.8cm},
  arr/.style={->, >=Stealth, thick, accent!70}
]

% Hexagonal placement of phases
\node[phase] (biz)  at (0,    3.5) {Business\\Understanding};
\node[phase] (data) at (3.2,  2.0) {Data\\Understanding};
\node[phase] (prep) at (3.2, -0.5) {Data\\Preparation};
\node[phase] (mod)  at (0,   -2.0) {Modeling\\(Agent Design)};
\node[phase] (eval) at (-3.2,-0.5) {Evaluation};
\node[phase] (dep)  at (-3.2, 2.0) {Deployment};

% Center label
\node[draw=accent!40, fill=yellow!10, rounded corners=4pt,
      font=\footnotesize\bfseries, align=center] at (0, 0.7)
  {Iterative\\CRISP-DM\\Cycle};

% Arrows (cycle)
\draw[arr] (biz)  -- (data);
\draw[arr] (data) -- (prep);
\draw[arr] (prep) -- (mod);
\draw[arr] (mod)  -- (eval);
\draw[arr] (eval) -- (dep);
\draw[arr] (dep)  -- (biz);

% Annotations
\node[annot, right=0.15cm of data]
  {DOM state, screenshots,\\network traces};
\node[annot, right=0.15cm of prep]
  {URL dedup, output shaping,\\context structuring};
\node[annot, below=0.2cm of mod]
  {Agent roles, tools,\\prompt policy};
\node[annot, left=0.15cm of eval]
  {Manual live-match\\campaigns + auto checks};
\node[annot, left=0.15cm of dep]
  {Azure Container Apps\\Jobs + Service Bus};
\node[annot, above=0.2cm of biz]
  {Produce evidence for\\takedown actions};

\end{tikzpicture}
\caption{CRISP-DM methodology adapted for agentic browser evidence collection.}
\label{fig:crispdm}
\end{figure}

\section{Project evolution timeline}

The project progressed through two main stages as shown in Figure~\ref{fig:timeline}.

\begin{figure}[H]
\centering
\begin{tikzpicture}[
  milestone/.style={draw=accent, fill=blue!8, rounded corners=4pt,
                    minimum width=3.8cm, minimum height=1.4cm,
                    text centered, font=\footnotesize, align=center},
  event/.style={font=\tiny, align=center, text width=3.0cm, text=gray!80}
]

% Main horizontal axis
\draw[very thick, accent!60] (-0.5,0) -- (13.5,0);
\foreach \x in {0,3.5,7,10.5,13}
  \draw[thick, accent!60] (\x,-0.15) -- (\x,0.15);

% Stage 1
\node[milestone, fill=orange!10, draw=orange!60] at (1.75, 1.5)
  {\textbf{n8n Phase}\\Rapid feasibility\\validation};
\draw[->, thick, orange!70] (1.75,0.9) -- (1.75,0.15);

\node[event] at (0,  -0.8) {Internship\\start};
\node[event] at (3.5,-0.8) {n8n\\prototype};

% Stage 2
\node[milestone, fill=blue!10, draw=blue!60] at (5.25, 1.5)
  {\textbf{Transition}\\Lessons codified;\\new platform design};
\draw[->, thick, blue!70] (5.25,0.9) -- (5.25,0.15);
\node[event] at (7,-0.8) {Python\\platform};

% Stage 3
\node[milestone, fill=green!10, draw=green!60] at (8.75, 1.5)
  {\textbf{Python Platform}\\FastAPI, agents,\\tools, observability};
\draw[->, thick, green!70] (8.75,0.9) -- (8.75,0.15);
\node[event] at (10.5,-0.8) {Evaluation\\campaigns};

% Stage 4
\node[milestone, fill=purple!10, draw=purple!60] at (11.75, 1.5)
  {\textbf{Hardening}\\Cloud packaging,\\deployment design};
\draw[->, thick, purple!70] (11.75,0.9) -- (11.75,0.15);
\node[event] at (13,-0.8) {Azure\\target design};

\end{tikzpicture}
\caption{Project evolution timeline: n8n phase to coded platform to cloud-ready design.}
\label{fig:timeline}
\end{figure}

\section{CRISP-DM applied to this project}

Each phase was adapted to the constraints of agentic browser automation on adversarial
targets. Business understanding produced the core success criteria: classify every URL
reliably, extract at least one validated stream URL per hosting page, and make every
run fully auditable. Data understanding was conducted through manual observation of
live pages — DOM output, screenshots, network traces, and frame trees — since no static
labelled dataset exists. Data preparation covers URL normalization, HTML noise reduction
(stripping scripts, ads, and tracking tags before LLM context building), and network
log filtering to isolate m3u8/mpd candidates from unrelated traffic. Modeling refers to
agent role design, prompt iteration (covered in Chapter~5), and model routing: lightweight
models handle cheap decisions (classification, orchestration) while stronger models run
extraction loops. Evaluation ran continuously through manual live-match campaigns and
automated regression checks. Deployment follows an event-driven cloud model: FastAPI
and Next.js as always-on services, Puppeteer extraction as Azure Container Apps Jobs
triggered by Azure Service Bus, and Cloudinary for stateless screenshot storage.

\section{Functional and non-functional requirements}

Table~\ref{tab:fr-summary} summarizes the ten functional requirements derived from the
operational context. All critical requirements were implemented in the final platform.

\begin{table}[H]
\centering
\begin{tabularx}{\textwidth}{>{\raggedright\arraybackslash\bfseries}p{1.2cm}
                              >{\raggedright\arraybackslash}X
                              >{\raggedright\arraybackslash}p{2.5cm}}
\toprule
\rowcolor{accent!8}
ID & Requirement summary & Priority \\
\midrule
FR-1 & URL ingestion and run lifecycle management & Critical \\
FR-2 & Page-type classification with DOM + screenshot & Critical \\
FR-3 & Specialist routing by page type & Critical \\
FR-4 & Tab-aware landing-page candidate discovery & High \\
FR-5 & Multi-server hosting extraction loop & Critical \\
FR-6 & Embedded player activation and stream capture & Critical \\
FR-7 & Hosting-to-embedded fallback & High \\
FR-8 & Screenshot capture and Cloudinary storage & High \\
FR-9 & WHOIS / IP infrastructure enrichment & Medium \\
FR-10 & Real-time operator console + SSE tracing & High \\
\bottomrule
\end{tabularx}
\caption{Functional requirements summary with priority.}
\label{tab:fr-summary}
\end{table}

Table~\ref{tab:nfr-summary} lists the non-functional requirements and the architectural
decisions that address each one.

\begin{table}[H]
\centering
\begin{tabularx}{\textwidth}{>{\raggedright\arraybackslash\bfseries}p{1.4cm}
                              >{\raggedright\arraybackslash}X
                              >{\raggedright\arraybackslash}p{3cm}}
\toprule
\rowcolor{accent!8}
ID & Non-functional requirement & How it is addressed \\
\midrule
NFR-1 & Maintainability & Versioned prompt files; modular tool registry \\
NFR-2 & Observability & Structured event log; SSE console; cost dashboard \\
NFR-3 & Cost control & Model routing; provider-side caching; token budgets \\
NFR-4 & Reliability & Tool error contracts; popup handler; retry logic \\
NFR-5 & Scalability & Azure Container Apps Jobs; Service Bus queue \\
NFR-6 & Portability & Docker Compose; configurable browser profiles \\
NFR-7 & Evidence traceability & Run-linked artifact chain: event $\to$ screenshot $\to$ URL \\
\bottomrule
\end{tabularx}
\caption{Non-functional requirements and implementation approach.}
\label{tab:nfr-summary}
\end{table}

\section{Design constraints}

Five constraints bounded the solution space throughout: (1) target-site layouts and
anti-bot behavior can change without notice between broadcast windows; (2) the player
can be nested behind multiple iframe hops from the top-level page; (3) popups, overlays,
and fake play controls can divert scripted extraction; (4) final stream manifests may
only be generated after a runtime token exchange triggered by player activation; and
(5) inference cost scales with exploration depth, number of hosting pages, and servers
tested per page — requiring active caching and model routing to remain viable at scale.

\section{Prompting methodology evolution: intent-based to ReAct-style loops}

The prompting strategy evolved in two stages. Figure~\ref{fig:prompt-evolution} compares
the two approaches side by side.

\begin{figure}[H]
\centering
\begin{tikzpicture}[node distance=0.55cm and 2.4cm]

% ==== LEFT: Intent-based ====
\node[flowbox, minimum width=5cm, fill=orange!12, draw=orange!60] (i1)
  at (0, 5.5) {\textbf{Intent-Based Prompt (early)}};
\node[toolbox, minimum width=4.8cm, below=0.5cm of i1] (i2)
  {System role + high-level goal};
\node[toolbox, minimum width=4.8cm, below=0.45cm of i2] (i3)
  {List of tools available};
\node[toolbox, minimum width=4.8cm, below=0.45cm of i3] (i4)
  {Expected JSON output schema};
\node[toolbox, minimum width=4.8cm, below=0.45cm of i4] (i5)
  {LLM calls tools in one pass};
\node[toolbox, minimum width=4.8cm, fill=red!10, draw=red!50, below=0.45cm of i5] (i6)
  {Returns answer — stops early\\if page is ambiguous};

\draw[grayarrow] (i1)--(i2); \draw[grayarrow] (i2)--(i3);
\draw[grayarrow] (i3)--(i4); \draw[grayarrow] (i4)--(i5);
\draw[grayarrow] (i5)--(i6);

% ==== RIGHT: ReAct ====
\node[flowbox, minimum width=5cm, fill=blue!12, draw=blue!60] (r1)
  at (7.5, 5.5) {\textbf{ReAct-Style Prompt (current)}};
\node[toolbox, minimum width=4.8cm, below=0.5cm of r1] (r2)
  {System role + final objective};
\node[toolbox, minimum width=4.8cm, below=0.45cm of r2] (r3)
  {Explicit loop instruction:\\Observe $\to$ Reason $\to$ Act};
\node[toolbox, minimum width=4.8cm, below=0.45cm of r3] (r4)
  {Evidence-grounded stopping\\criteria};
\node[toolbox, minimum width=4.8cm, below=0.45cm of r4] (r5)
  {LLM iterates: tool call $\to$\\re-observe $\to$ next decision};
\node[toolbox, minimum width=4.8cm, fill=green!10, draw=green!50,
      below=0.45cm of r5] (r6)
  {Stops when evidence found\\or budget hit};
\draw[->, >=Stealth, thick, blue!50, bend right=35] (r6.east) to (r5.east);

\draw[myarrow] (r1)--(r2); \draw[myarrow] (r2)--(r3);
\draw[myarrow] (r3)--(r4); \draw[myarrow] (r4)--(r5);
\draw[myarrow] (r5)--(r6);

% Divider
\draw[dashed, gray!50, thick] (3.75,-0.3) -- (3.75,6.0);

\end{tikzpicture}
\caption{Prompt strategy evolution: intent-based linear flow versus ReAct-style
iterative loop.}
\label{fig:prompt-evolution}
\end{figure}

The practical difference was not subtle. The intent-based prompt would often declare a run complete after inspecting the first server option, whether or not a stream had actually appeared in network traffic. The ReAct loop couldn't do that — the stopping criteria required explicit evidence. This also made debugging much easier, since each step in the trace corresponded to a specific reasoning decision rather than one large opaque LLM call.

The design principles that fell out of this analysis were: classify before deep extraction; use specialist agents with explicit fallback paths; ground every conclusion in browser evidence; validate state changes with screenshots; control cost through caching and model routing; and treat anti-bot configuration as a maintenance process, not a one-time setup.

The requirements, constraints, and methodology defined in this chapter establish the
formal scope of the system. Chapter~4 translates these inputs into a concrete
multi-agent architecture, documenting the component structure, agent roles, tool
allocation decisions, and deployment model.
